\documentclass[9pt,twocolumn,twoside]{article} %210 mm × 297 mm
\usepackage[utf8]{inputenc}
\usepackage[spanish, es-tabla]{babel}
\spanishdecimal{.}
\usepackage[a4paper,left=1.5cm, right=1.5cm, top=2cm, bottom=2cm]{geometry}
\usepackage{graphicx}
\usepackage{blindtext}
\usepackage{multirow}
\usepackage{mhchem}
\usepackage{url}

\newcommand{\titulo}{Sustituye esto por el título del artículo}
\newcommand{\nombre}{Sustituye esto por tu nombre}

\newcommand{\resumen}{Aquí hay que escribir un resumen del artículo. Debe de ser un resumen preciso sin ser muy extenso. Sin embargo, debe dejar claro de qué trata el trabajo.}


\usepackage{fancyhdr}
\pagestyle{fancy}
\fancyhead{}
\fancyfoot{}
\fancyhead[LO]{\small\nombre}
\fancyhead[RE]{\small\titulo}
\fancyfoot[RO,LE] {\thepage}

\setcounter{page}{1} %Número de la página, la editorial lo cambiará

\begin{document}

\twocolumn[
\begin{@twocolumnfalse}
\begin{tabular}{p{5cm} p{12cm}}
\multirow{2}{*}{\includegraphics[scale=0.2]{LOGO_QUIMICA_3.png}}&{{\Large\bf\titulo}}\\
\vspace{0.2cm} & \vspace{0.2cm}\\
 &{\large\nombre}\\
\vspace{0.1cm} & \vspace{0.1cm}\\ 
 & {\small\it Departamento de Química. Universidad Autónoma Metropolitana Unidad Iztapalapa. San Rafael Atlixco 186. Colonia Vicentina. Iztapalapa. C.P. 09340. Ciudad de México, México.}\\
{\tiny Fecha de recepción: No escribir nada} &\\
{\tiny Fecha de aceptación: No escribir nada} & {\small\resumen}\\
\end{tabular}
\end{@twocolumnfalse}
\vspace{0.5cm}
]

\section{Introducción}
Escribe los antecedentes y la motivación del trabajo, sobre todo debes de escribir el estado del arte de tu tema y plantear el problema que surge de este análisis.

\section{Ecuaciones y fórmulas}
El uso de ecuaciones en LaTex es diferente al que se usa en un procesador de texto. Sin embargo, LaTex es versátil y muy potente con el uso de símbolos. A continuación se tienen algunos ejemplos:

\begin{equation}
\left(p+\frac{n^2 a}{V^2}\right)\left(V-nb\right)=nRT.\label{ec:gases_1}
\end{equation}

Otra forma de escribir una ecuación dentro de un texto es de la siguiente manera $p=\left(\partial E/\partial V\right)_T$. Para más símbolos dentro de ecuaciones, recurre al sitio \url{https://oeis.org/wiki/List_of_LaTeX_mathematical_symbols}.

La ecuación (\ref{ec:gases_1}) es importante para la descripción de gases no ideales. Para una reacción química podemos escribir fórmulas químicas \ce{H2O}, \ce{SO4^2-}, \ce{BaSO4 v}.
Y una reacción química puede ser escrita de la siguiente manera

\begin{center}
\ce{CH4 ->[O2][\Delta] CO2 + H2O}
\end{center}

En la siguiente dirección encontrarás toda la información relacionada con el paquete mhchem: \url{https://www.ctan.org/pkg/mhchem}

\section{Referencias}
Un artículo científico se sustenta en trabajo previo y por lo tanto el uso de referencias es muy importante. Las citas en nuestra revista se harán de acuerdo al uso de bibtex integrado en LaTex. En el archivo referencias se presentan algunos ejemplos típicos que se usa en nuestra literatura. Para citar alguna referencia que se encuentra en nuestro archivos referencia solamente hay que usar $\setminus$\texttt{cite\{Referencia\}}, por ejemplo una cita al profesor R. Sagar \cite{Sagar_1992} o la cita a un libro \cite{Szabo}. Con bibtex no tenemos que preocuparnos por la numeración de las referencias. Por lo tanto solamente tenemos que preocuparnos por escribir las referencias en el formato mostrado en el archivo \texttt{referencias.bib}.

\section{Figuras}
Las figuras deben ser insertadas en LaTex en formato png o eps. En la Figura \ref{fig:prueba} se presenta un ejemplo.
\begin{figure}[h]
    \centering
    \includegraphics[width=8cm]{Tlecaxitl_logo.png}
    \caption{Aquí debemos de poner la descipción de la figura.}
    \label{fig:prueba}
\end{figure}
En este caso se utiliza el ancho máximo permitido en nuestra revista (8 cm).
\section{Tablas}
El reporte de resultados a través de tablas es muy importante para un artículo científico. Sugerimos los siguientes formatos de tablas. En la Tabla \ref{tab:uno} se presenta un formato donde un número se escribe como una ecuación matemática ya que usa el símbolo \$ para delimitar el número 1.61. En este caso no hay diferencia alguna si no se escribe el número entre este símbolo o no. El ejemplo  reporta una tabla con tres columnas, definidas cuando se escribe el comando tabular. Las primeras dos columnas están justificadas hacia la izquierda y la tercera se encuentra centrada. Entre las tres columnas se declara una raya vertical, todo esto se hace con el comando $\setminus$\texttt{begin\{tabular\}\{l|l|c\} (left, left y center)}.
En ocasiones es necesario hacer referencia a resultados ya publicados y que se encuentran en el archivos. En la Tabla \ref{tab:uno} se presenta un ejemplo de este tipo. 

\begin{table}[h]
\caption{Este es un ejemplo de tabla en nuestra revista.}
\label{tab:uno}
\begin{center}
\begin{tabular}{l|l|c}
\hline\hline
Nombre      &   Apellidos       &   Altura (m)\\
\hline
Enrique     &   Bonilla Cerros  &   $1.61$ \\
Carolina    &   Paredes Heredia &   1.71$^a$\\
\hline\hline
\end{tabular}
\end{center}
$^a$ Ref. \cite{1973_MEP_1}.
\end{table}

\section{Conclusiones}
Para generar documentos LaTex existe un gran número de ejemplos que pueden ser encontrados en muchas páginas web o en YouTube. Lo más importante es que este documento te debe servir para armar todo tu artículo, por favor no incluyas \texttt{package} adicionales.

\section*{Agradecimientos}
En esta sección escribe los agradecimientos a las personas o instituciones que te permitieron desaroolar este trabajo.

\bibliographystyle{unsrt}
\bibliography{referencias}
\end{document}
